\clearpage
\fancyhead[L]{\small VON LINDENAU. STERNWARTE SEEBERG.}
\hypertarget{native-503-lower}{}\label{native-503-lower}
\newcommand{\SeebergResidualMark}{\hbox{\vrule width0pt height3.272222bp depth0pt\special{pdf:bcontent}\special{pdf:code q 0 g 2.927778 1.377778 m 2.927778 1.205556 l 2.755556 1.033333 l 2.755556 0.172222 l 2.927778 0.000000 l 3.100000 0.172222 l 3.100000 1.377778 l h 0.172222 3.272222 m 0.000000 3.100000 l 0.000000 2.927778 l 0.172222 2.755556 l 1.205556 2.755556 l 1.377778 2.927778 l 2.066667 2.927778 l 2.238889 2.755556 l 2.755556 2.755556 l 2.927778 2.583333 l 2.927778 2.411111 l 3.100000 2.411111 l 3.100000 2.583333 l 3.444444 2.927778 l 6.027778 2.927778 l 6.200000 3.100000 l 6.027778 3.272222 l 3.961111 3.272222 l 3.788889 3.100000 l 3.616667 3.100000 l 3.444444 3.272222 l 2.583333 3.272222 l 2.411111 3.100000 l 0.516667 3.100000 l 0.344444 3.272222 l h f* Q}\special{pdf:econtent}\kern6.200000bp}}
\newcommand{\SeebergMoon}{\text{\mdseries\AstroSymbols ☽}}

\begin{center}
\small Göttingische gelehrte Anzeigen. Stück 106. Seite 1049..1052. 1808 Juli 2.\\[3pt]
\rule{.72\linewidth}{.3pt}
\end{center}

Ueber den gegenwärtigen Zustand der berühmten Sternwarte \emph{Seeberg} ist uns kürzlich von Hrn. \textsc{von Lindenau} eine authentische Nachricht in einem besondern Aufsatze zugesandt worden, woraus ein Auszug den Verehrern der Himmelskunde um so willkommner sein wird, da hierüber seit einiger Zeit aus trüben Quellen manche ungegründete Nachrichten ins Publicum gebracht waren. Die Kriegsunruhen, verbunden mit einigen am Dache der Sternwarte erforderlich gewordenen Reparaturen, hatten im Jahre 1806 die Abnahme der Instrumente nothwendig gemacht: mancherlei Umstände, deren Aufzählung nicht hieher gehört, erlaubten erst zu Anfang dieses Jahres, wegen Wiederaufstellung derselben die nöthigen Anstalten zu treffen. Hr. \textsc{von Lindenau}, welcher schon einmal, in Abwesenheit des Hrn. \textsc{von Zach}, die Aufsicht über die Sternwarte geführt hatte, wurde von des Herzogs von Sachsen-Gotha Durchl. mit jenem Geschäfte beauftragt, und nahm am 9. April auf der Sternwarte seine Wohnung. Erst am 20. April konnte, der widrigen Witterung wegen, mit der Wiederaufstellung der Instrumente der Anfang gemacht werden: inzwischen wurde nachher dies Geschäft von mehreren hellen Nächten so gut begünstigt, dass Ende Aprils Passage-Instrument, Quadrant und Regulator in beobachtungsfertigem Zustande war. Wie vollkommen dem Hrn. \textsc{von Lindenau} die Berichtigung der Instrumente sogleich gelungen ist, beurkundet eine schöne Reihe zahlreicher, vom 29. April bis 9. Mai angestellter und uns mitgetheilter, Beobachtungen, von denen wir hier, des beschränkten Raumes dieser Blätter wegen, nur einige wiedergeben können:

\begin{center}
\emph{Beobachtungen der Sonne.}

\smallskip
{\small\setlength{\tabcolsep}{3pt}
\begin{tabular}{r|rrr|rrr|r}
1808 & \multicolumn{3}{c|}{\shortstack{Mittlere Zeit\\in Seeberg}} & \multicolumn{3}{c|}{\shortstack{Gerade Aufsteigung\\der Sonne}} & \shortstack{Fehler der\\Tafeln} \\
\hline
April 30 & $23^{\mathrm{u}}$ & $57^{\mathrm{m}}$ & $4^{\mathrm{s}}16$ & $37^{\circ}$ & $29'$ & $4''2$ & $+3''8$ \\
Mai 3 & 23 & 56 & 42.06 & 40 & 20 & 57.4 & $+1.0$ \\
5 & 23 & 56 & 29.87 & 42 & 16 & 9.9 & $+2.5$ \\
6 & 23 & 56 & 24.89 & 43 & 14 & 3.3 & $-2.0$ \\
7 & 23 & 56 & 20.06 & 44 & 11 & 59.0 & $-0.5$ \\
9 & 23 & 56 & 11.73 & 46 & 8 & 15.5 & $\SeebergResidualMark\,2.6$ \\
\hline
\multicolumn{7}{r|}{Fehler der Tafeln im Mittel} & $+1.8$ \\
\end{tabular}}
\end{center}

\clearpage
\SourceLeaf{504}
\begin{center}
\emph{Beobachtungen des Mondes.}

\smallskip
{\small\setlength{\tabcolsep}{3pt}
\begin{tabular}{r|rrr|rrr|r}
1808 & \multicolumn{3}{c|}{\shortstack{Mittlere Zeit\\in Seeberg}} & \multicolumn{3}{c|}{\shortstack{Ger. Aufsteig. des\\westl. $\SeebergMoon$randes}} & \shortstack{Fehler der\\Tafeln.} \\
\hline
Mai 2 & $5^{\mathrm{u}}$ & $20^{\mathrm{m}}$ & $37^{\mathrm{s}}63$ & $120^{\circ}$ & $34'$ & $0''3$ & $+2''5$ \\
3 & 6 & 10 & 48.31 & 134 & 7 & 50.4 & $+0.6$ \\
4 & 7 & 0 & 59.58 & 147 & 41 & 52.8 & $+5.7$ \\
5 & 7 & 51 & 44.20 & 161 & 24 & 15.9 & $+6.2$ \\
6 & 8 & 43 & 44.35 & 175 & 25 & 30.6 & $+3.1$ \\
7 & 9 & 37 & 44.89 & 189 & 57 & 4.3 & $+2.3$ \\
8 & 10 & 13 & 19.15 & 205 & 7 & 5.8 & $+6.7$ \\
\hline
\multicolumn{7}{r|}{Mittlerer Fehler der Tafeln} & $+3''9$ \\
\end{tabular}}
\end{center}

Die Tafeln, womit die Beobachtungen verglichen wurden, sind die neuen Sonnentafeln des Hrn. \textsc{von Zach}, und die Mondstafeln von \textsc{Bürg}.

Ausserdem theilte uns Hr. \textsc{von Lindenau} noch 8 beobachtete Oerter des \emph{Uranus} mit, 11 Oerter des \emph{Saturn}, 6 Oerter des \emph{Jupiter}, 9 Oerter der Venus, 8 Oerter des \emph{Merkur}: wir begnügen uns, hier nur die von Hrn. \textsc{von Lindenau} daraus gezogenen Endresultate anzuführen. Der mittlere Fehler der \textsc{de Lambre}'schen Uranustafeln für diese Epoche war $-0''8$ in der Länge, und $-16''$ in der Breite, der mittlere Fehler der Saturnstafeln in der Länge $+9''1$, in der Breite $-4''3$; die Opposition des Saturn 1808, Mai 9, $8^{\mathrm{u}}25^{\mathrm{m}}47^{\mathrm{s}}2$ mittl. Zeit in Seeberg; der mittlere Fehler der Jupiterstafeln in der Länge $-1''6$, in der Breite $-1''9$. Die Beobachtungen der Venus und des Mercur hat Hr. \textsc{von Lindenau} nicht verglichen, weil er die Tafeln dieser Planeten dem heutigen Zustande der Astronomie nicht mehr angemessen fand. Seit einem Jahre beschäftigt er sich selbst mit Sammlung von Materialien zu neuen Venustafeln, wozu er schon mehr als 200 gute Beobachtungen zusammengebracht hat; hiedurch, und mit Benutzung der von \textsc{la Place} entwickelten Störungsgleichungen, hofft er bald etwas Vollkommneres, als die bisherigen Tafeln, liefern zu können.

Als Hauptgegenstand seiner practischen Beschäftigungen während seines Aufenthalts auf der Seeberger Sternwarte hat sich Hr. \textsc{von Lindenau} die Bestimmung der Refractionen und der jährlichen Parallaxe der Fixsterne in gerader Aufsteigung vorgesetzt. Wenn es gegründet ist, dass $\alpha$ Lyrae eine Declinations-Parallaxe von $4\ldots5''$ zeigt, so muss die Parallaxe in gerader Aufsteigung $7''$ betragen, also der Unterschied des positiven und negativen Maximum fast auf eine Zeitsecunde steigen. Wird also ein solcher Stern mit mehreren in der Nähe befindlichen kleinen verglichen, so darf man hoffen, dass eine solche Differenz dem so vortrefflichen 8füssigen \textsc{Ramsden}schen Passage-Instrument nicht entgehen wird. Freilich wird die Vergleichung mit kleinen Sternen im positiven und negativen Maximum bedeutenden \emph{practischen} Schwierigkeiten unterworfen sein.

Die Beobachtungen der Refractionen hat Hr. \textsc{von Lindenau} bereits angefangen: obgleich es noch zu früh ist, ein bestimmtes Resultat schon jetzt daraus zu ziehen, so vereinigen sie sich doch alle dahin, bei einem Barometerstande von 28 Zoll und einer Temperatur von $10^{\circ}$ Reaumür eine Horizontalrefraction von $37\ldots40'$ zu geben. (Wir bemerken bei dieser Gelegenheit, dass auch Hr. Inspector \textsc{Bessel} aus einer sehr sorgfältigen Discussion der \textsc{Bradley}'schen Beobachtungen die Nothwendigkeit, die gewöhnliche Angabe der Horizontalrefraction beträchtlich vergrössern zu müssen, geschlossen hat). — Ueberhaupt wird Hr. \textsc{von Lindenau} nach Jahresfrist die Resultate seines Aufenthalts auf der Seeberger Sternwarte den Astronomen umständlich vorlegen.

Das ganze, zur eigentlichen Sternwarte gehörige, Gebäude ist jetzt völlig wieder hergestellt, und an der Wiedererbauung des Wohngebäudes wird gearbeitet, so dass es hoffentlich ebenfalls noch im Laufe dieses Sommers vollendet sein wird.
